\lecture{8.2}{Sequential Identification Assumptions}
\label{lec:8-2}

\subsection{Sequential Randomized Trials}

The cleanest way to learn the mean counterfactual outcome $\E\{Y(\abar)\}$ is
to run a trial that randomizes treatment at each time point. In a sequentially
randomized trial, treatment $A_k$ is assigned by a known randomization mechanism
after observing the history available at time $k$.

For a fixed regime $\abar=(a_0,\ldots,a_K)$, randomization implies that treatment
at each time is independent of the relevant counterfactual outcome within the
history strata used by the trial:
\[
  Y(\abar) \ci A_k \given \Lbar_k,\Abar_{k-1}
  \qquad \text{for } k=0,\ldots,K.
\]
At $k=0$, the conditioning set is just $L_0$.

Plainly, among people with the same history up to time $k$, the next treatment
assignment is as if randomized with respect to what would happen under the
regime $\abar$. This is the longitudinal analogue of conditional randomization
for a point treatment.

\subsection{Observational Longitudinal Studies}

In observational longitudinal studies, treatment is not assigned by the analyst.
Patients and clinicians make decisions using information that evolves over time.
Some of this information is measured and included in $\Lbar_k$; some may be
unmeasured.

The inferential challenge is that the observed treatment process may be
informative about future outcomes. For example, a patient may discontinue
chemotherapy because disease is progressing, or a smoker may quit because early
respiratory symptoms have appeared. If those health changes are also associated
with the final outcome, treated and untreated histories are not automatically
comparable.

\subsection{Time-Varying Confounding}

\begin{definition}[Time-varying confounder]
A time-varying confounder is a covariate $L_k$ that predicts later treatment
$A_k$ or $A_{k+1}$ and also predicts the outcome, while possibly being affected
by earlier treatment.
\end{definition}

This dual role is what makes longitudinal confounding different from ordinary
baseline confounding. A later covariate can be necessary for controlling
confounding of later treatment, yet inappropriate to condition on naively when
estimating effects of earlier treatment because it lies on a causal pathway from
earlier treatment to the outcome.

\begin{figure}[H]
\centering
\begin{tikzpicture}[
  node distance=1.2cm,
  box/.style={draw, rounded corners=2pt, minimum width=1.0cm, minimum height=.72cm},
  arr/.style={-{Latex[length=2.3mm]}, thick}
]
\node[box] (L0) {$L_0$};
\node[box, right=of L0] (A0) {$A_0$};
\node[box, right=of A0] (L1) {$L_1$};
\node[box, right=of L1] (A1) {$A_1$};
\node[box, right=of A1] (Y) {$Y$};
\draw[arr] (L0) -- (A0);
\draw[arr] (L0) to[bend left=25] (L1);
\draw[arr] (L0) to[bend left=35] (Y);
\draw[arr] (A0) -- (L1);
\draw[arr] (A0) to[bend left=25] (Y);
\draw[arr] (L1) -- (A1);
\draw[arr] (L1) to[bend left=25] (Y);
\draw[arr] (A1) -- (Y);
\end{tikzpicture}
\caption{Treatment-confounder feedback: $L_1$ is affected by $A_0$ and also confounds the effect of $A_1$ on $Y$.}
\label{fig:ch8-time-varying-confounding}
\end{figure}

\subsection{Sequential Randomization}

\begin{assumptionbox}
For a fixed static regime $\abar=(a_0,\ldots,a_K)$, the sequential
randomization assumption is
\[
  Y(\abar) \ci A_k \given \Lbar_k,\Abar_{k-1}=\abar_{k-1}
  \qquad k=0,\ldots,K.
\]
Equivalently, among individuals whose observed treatment history has followed
the regime through time $k-1$, the next treatment decision is independent of the
counterfactual outcome under the full regime after conditioning on the measured
covariate history.
\end{assumptionbox}

The restriction $\Abar_{k-1}=\abar_{k-1}$ appears because, when identifying
$Y(\abar)$, the relevant risk set at time $k$ consists of individuals who have
remained compatible with the regime up to that time. If a person has already
deviated from the regime, later treatment decisions for that person do not
directly reveal outcomes under the full regime $\abar$ without further modeling.

\subsection{Sequential Exchangeability}

Sequential exchangeability is another name for the same identifying condition.
It emphasizes comparability: after conditioning on the measured history, those
who receive $A_k=a_k$ and those who do not are exchangeable with respect to the
future counterfactual outcome under $\abar$.

For two treatment times and regime $(1,1)$, the assumption is
\[
  Y(1,1) \ci A_0 \given L_0,
\]
and
\[
  Y(1,1) \ci A_1 \given L_0,A_0=1,L_1.
\]
The first line says baseline treatment is unconfounded within baseline covariate
strata. The second says the second treatment is unconfounded within strata of
the history observed among people treated at baseline.

This assumption is not empirically verifiable. The observed data can show
whether treatment depends on measured covariates, but cannot show whether
treatment depends on unmeasured determinants of the counterfactual outcome.

\subsection{Longitudinal Positivity}

\begin{assumptionbox}
For every history with positive probability under the regime of interest,
the probability of receiving the regime-compatible treatment must be positive:
\[
  \Prob(A_k=a_k \given \Lbar_k=\lbar_k,\Abar_{k-1}=\abar_{k-1})>0
  \qquad k=0,\ldots,K.
\]
\end{assumptionbox}

For the always-treated regime $(1,\ldots,1)$ this becomes
\[
  \Prob(A_0=1\given L_0=l_0)>0,
\]
\[
  \Prob(A_1=1\given \Lbar_1=\lbar_1,A_0=1)>0,
\]
and so on through time $K$.

Longitudinal positivity is more fragile than single-timepoint positivity because
probabilities multiply over time. Even if every time-specific probability is
moderate, the probability of observing the entire treatment sequence can be
small:
\[
  \prod_{k=0}^K
  \Prob(A_k=a_k\given \Lbar_k,\Abar_{k-1}=\abar_{k-1}).
\]
Small products lead to unstable inverse probability weights and can make some
static regimes impossible to learn in a given population.

\begin{keyidea}
Identification of a longitudinal regime requires local exchangeability and local
positivity at every treatment decision. A failure at one time point can break
identification of the entire treatment sequence.
\end{keyidea}
